\documentclass[a4paper,12pt]{article}
\usepackage[T2A]{fontenc}
\usepackage[utf8]{inputenc}
\usepackage[russian]{babel}
\usepackage{amsmath,amssymb}
\usepackage{siunitx}
\usepackage{booktabs}
\usepackage{geometry}
\geometry{margin=25mm}

\begin{document}
	
	\subsubsection*{Экспериментальные данные}
	\subsubsection*{Опыт 1}
	\begin{table}[h]
		\centering
		\begin{tabular}{lcc}
			\toprule
			Раствор & Окраска лакмуса & Вывод о pH \\
			\midrule
			$\mathrm{NaCl}$ & фиолетовая & нейтральная среда (pH$\approx$7) \\
			$\mathrm{Na_2CO_3}$ & синяя & щелочная среда (pH>7) \\
			$\mathrm{AlCl_3}$ & красная & кислая среда (pH<7) \\
			Дистиллированная вода & фиолетовая & нейтральная (pH$\approx$7) \\
			\bottomrule
		\end{tabular}
	\end{table}
	
	\subsubsection*{Расчёт для $\mathrm{Na_2CO_3}$}
	\[
	\mathrm{CO_3^{2-}+H_2O \rightleftharpoons HCO_3^-+OH^-}
	\]
	\[
	K_b=\frac{K_w}{K_{a2}}=\frac{10^{-14}}{4.7\cdot10^{-11}}\approx 2.13\cdot10^{-4}.
	\]
	\[
	[\mathrm{OH^-}] \approx \sqrt{K_b c}\approx 1.46\cdot10^{-2}.
	\]
	\[
	\mathrm{pOH}=1.84,\qquad \mathrm{pH}\approx 12.16.
	\]
	Степень гидролиза:
	\[
	h=\frac{[\mathrm{OH^-}]}{c}\approx 0.0146.
	\]
	
	\subsubsection*{Расчёт для $\mathrm{AlCl_3}$}
	\[
	\mathrm{[Al(H_2O)_6]^{3+} \rightleftharpoons [Al(H_2O)_5(OH)]^{2+}+H^+}
	\]
	\[
	[\mathrm{H^+}] \approx \sqrt{K_a c}\approx 3.16\cdot10^{-3},
	\]
	\[
	\mathrm{pH}\approx 2.5,
	\qquad h\approx 0.00316.
	\]
	
	\subsubsection*{Основные уравнения реакций}
	\begin{itemize}
		\item Гидролиз карбоната:
		\[
		\mathrm{CO_3^{2-}+H_2O \rightleftharpoons HCO_3^- + OH^-}
		\]
		\item Гидролиз алюминия:
		\[
		\mathrm{[Al(H_2O)_6]^{3+} \rightleftharpoons [Al(H_2O)_5(OH)]^{2+}+H^+}
		\]
		\item Гетерогенный гидролиз:
		\[
		\mathrm{Al^{3+}+3CO_3^{2-}+3H_2O \to Al(OH)_3 \downarrow +3HCO_3^-}
		\]
	\end{itemize}
	
	\subsubsection*{Выводы}
	\begin{enumerate}
		\item Соли сильных кислот и оснований ($\mathrm{NaCl}$) не подвергаются гидролизу.
		\item Соли сильных оснований и слабых кислот ($\mathrm{Na_2CO_3}$) образуют щелочную среду; рассчитано pH$\approx12.16$.
		\item Соли слабых оснований и сильных кислот ($\mathrm{AlCl_3}$) дают кислую реакцию среды (pH$\approx2.5$).
		\item Гетерогенный гидролиз сопровождается выпадением осадков ($\mathrm{Al(OH)_3}$, $\mathrm{Cu(OH)_2}$ и др.).
		\item Нагревание и разбавление усиливают гидролиз согласно принципу Ле Шателье.
	\end{enumerate}
	
	\subsubsection*{Опыт 2. Гетерогенный гидролиз при добавлении $\mathrm{Na_2CO_3}$}
	
	\textbf{Наблюдения:}
	При добавлении раствора $\mathrm{Na_2CO_3}$ к раствору $\mathrm{CuCl_2}$ наблюдается выпадение 
	голубого осадка $\mathrm{Cu(OH)_2}$. 
	При добавлении $\mathrm{Na_2CO_3}$ к раствору $\mathrm{AlCl_3}$ выпадает белый студенистый осадок $\mathrm{Al(OH)_3}$. 
	Оба осадка растворяются в избытке кислоты.
	
	\textbf{Объяснение:}
	Анион $\mathrm{CO_3^{2-}}$ подвергается гидролизу:
	\[
	\mathrm{CO_3^{2-} + H_2O \rightleftharpoons HCO_3^- + OH^-}
	\]
	Избыток $\mathrm{OH^-}$ приводит к образованию малорастворимых гидроксидов:
	\[
	\mathrm{Cu^{2+} + 2OH^- \rightarrow Cu(OH)_2 \downarrow}
	\]
	\[
	\mathrm{Al^{3+} + 3OH^- \rightarrow Al(OH)_3 \downarrow}
	\]
	
	$\mathrm{Al(OH)_3}$ является амфотерным и растворяется в избытке щёлочи:
	\[
	\mathrm{Al(OH)_3 + OH^- \rightarrow [Al(OH)_4]^-}
	\]
	
	\textbf{Вывод:} 
	Гидролиз карбонат-ионов приводит к образованию $\mathrm{OH^-}$, который вызывает 
	осаждение гидроксидов слабых оснований. Данный опыт подтверждает 
	гетерогенную природу гидролиза и амфотерность $\mathrm{Al(OH)_3}$.
	
	\subsubsection*{Опыт 3. Сравнение силы связанных кислот и оснований}
	
	\textbf{Наблюдения:}
	Используя кислотно-основные индикаторы (лакмус, фенолфталеин), сравнивали 
	поведение растворов солей, образованных сильными и слабыми электролитами. 
	Растворы солей сильного основания и слабой кислоты (например, ацетат натрия) 
	дают щелочную реакцию. Соли слабого основания и сильной кислоты 
	(например, $\mathrm{NH_4Cl}$ или $\mathrm{AlCl_3}$) дают кислую реакцию.
	
	\textbf{Объяснение:}
	Гидролиз аниона слабой кислоты (например, $\mathrm{CH_3COO^-}$):
	\[
	\mathrm{CH_3COO^- + H_2O \rightleftharpoons CH_3COOH + OH^-}
	\]
	Гидролиз катиона слабого основания:
	\[
	\mathrm{NH_4^+ + H_2O \rightleftharpoons NH_3 \cdot H_2O + H^+}
	\]
	
	\textbf{Вывод:}
	Характер среды раствора соли определяется силой кислоты и основания, образующих соль:  
	\begin{itemize}
		\item сильная кислота + сильное основание → гидролиз отсутствует;
		\item слабая кислота + сильное основание → раствор становится щелочным;
		\item сильная кислота + слабое основание → раствор становится кислым.
	\end{itemize}
	
	\subsubsection*{Опыт 4. Влияние температуры на гидролиз соли}
	
	\textbf{Наблюдения:}
	При нагревании раствора $\mathrm{CH_3COONa}$ с добавленным индикатором наблюдается 
	усиление окрашивания в щелочную сторону. При охлаждении интенсивность 
	щелочной окраски уменьшается.
	
	\textbf{Объяснение:}
	Гидролиз солей слабых кислот является эндотермическим:
	\[
	\mathrm{CH_3COO^- + H_2O \rightleftharpoons CH_3COOH + OH^-}
	\]
	Повышение температуры смещает равновесие вправо, увеличивая концентрацию $\mathrm{OH^-}$,  
	что приводит к повышению pH.
	
	\textbf{Вывод:}
	Гидролиз ацетат-ионов усиливается при нагревании, так как процесс эндотермический.
	Температура является важным фактором, влияющим на степень гидролиза.
	
	\subsubsection*{Опыт 5. Влияние разбавления на гидролиз}
	
	\textbf{Наблюдения:}
	При постепенном разбавлении раствора $\mathrm{SbCl_3}$ водой появляется мутность, 
	а затем выпадает белый осадок $\mathrm{Sb(OH)_3}$.  
	При добавлении небольшого количества $\mathrm{HCl}$ осадок растворяется.
	
	\textbf{Объяснение:}
	Катион $\mathrm{Sb^{3+}}$ гидролизуется:
	\[
	\mathrm{Sb^{3+} + 3H_2O \rightleftharpoons Sb(OH)_3 \downarrow + 3H^+}
	\]
	Разбавление уменьшает концентрацию $\mathrm{H^+}$ и, согласно принципу Ле Шателье, 
	гидролиз усиливается, что приводит к выпадению гидроксида.  
	Добавление кислоты сдвигает равновесие влево, растворяя осадок.
	
	\textbf{Вывод:}
	Разбавление усиливает гидролиз солей слабых оснований.  
	Чем меньше концентрация раствора, тем больше степень гидролиза.
\end{document}